% =====================================================================
%  PARTIE 4 — Bénéfices et domaines d'application
% =====================================================================
\chapter{Bénéfices et domaines d'application}
\label{ch:benefices}

\section{Les bénéfices de l'AIoT et du TinyML}

L'association de l'intelligence artificielle et de l'Internet des objets,
appelée AIoT, permet aux équipements connectés d'analyser les données qu'ils
collectent afin de reconnaître des situations ou de détecter des anomalies. Le
TinyML correspond plus précisément à l'exécution de modèles d'apprentissage
automatique sur des dispositifs disposant de ressources limitées, notamment des
microcontrôleurs. Cette approche déplace une partie de l'intelligence au plus
près des capteurs, plutôt que de transmettre systématiquement les données à un
serveur distant. Elle présente plusieurs avantages en matière de réactivité, de
confidentialité et de consommation énergétique~\cite{ray2022review}.

\subsection{Une meilleure réactivité et une moindre dépendance au réseau}

Dans une architecture reposant sur le cloud, les données doivent être
transmises à un serveur avant d'être analysées. Le délai de réponse dépend donc
non seulement du calcul, mais aussi de la connexion réseau et des échanges avec
le serveur. Avec le TinyML, l'analyse peut être réalisée directement dans
l'objet connecté, ce qui évite cet aller-retour pour les décisions prises
localement.

Cette organisation est particulièrement intéressante lorsqu'un système doit
réagir rapidement, par exemple pour reconnaître un événement sonore ou détecter
un comportement inhabituel d'une machine. Elle permet également de maintenir
certaines fonctions d'analyse lorsque la connexion Internet est indisponible, à
condition que ces fonctions aient été conçues pour fonctionner de manière
autonome.

Il serait toutefois incorrect de parler de \og latence zéro\fg{}.
L'acquisition des données, leur préparation et l'exécution du modèle
nécessitent toujours du temps. Le bénéfice correspond à une réduction de la
latence et de sa dépendance au réseau, et non à sa disparition.

\subsection{Une meilleure maîtrise de la confidentialité des données}

Les objets connectés peuvent collecter des informations sensibles, comme des
enregistrements audio, des données de mouvement ou des mesures physiologiques.
Lorsque ces données sont systématiquement envoyées vers un serveur, leur
traitement implique davantage de transferts et de stockage à distance.
L'inférence locale permet de limiter ces échanges en analysant les informations
directement sur l'appareil.

Par exemple, un dispositif de reconnaissance sonore peut analyser un signal
audio localement et ne transmettre que le résultat de la détection. Selon
l'architecture retenue, il devient ainsi possible de communiquer une alerte
plutôt que l'ensemble de l'enregistrement. Le TinyML contribue donc à réduire
la quantité de données brutes qui quittent l'objet connecté.

Cet avantage ne signifie pas que la confidentialité est automatiquement
garantie. Certains systèmes peuvent encore transmettre des données pour
l'entraînement, le suivi ou d'autres services. Il faut donc distinguer la
possibilité de traiter les données localement du choix réel de ne pas les
transmettre.

\subsection{Une réduction possible de la consommation énergétique}

L'autonomie constitue une contrainte importante pour les objets connectés
alimentés par batterie. Les communications sans fil consomment de l'énergie,
notamment lorsque l'appareil transmet fréquemment des mesures. En analysant les
données localement et en ne communiquant que les résultats utiles, le TinyML
peut réduire le recours à la transmission radio~\cite{warden2019tinyml}.

Cependant, le calcul local consomme lui aussi de l'énergie. Le gain final
dépend donc de plusieurs paramètres : la complexité du modèle, la fréquence des
analyses, la consommation des capteurs, les modes de veille et la fréquence des
communications. Une analyse locale bien dimensionnée peut améliorer
l'autonomie, mais cet avantage doit être mesuré sur le dispositif complet
plutôt que supposé à partir du seul modèle d'IA~\cite{banbury2021mlperftiny}.

\begin{table}[H]
\centering
\small
\begin{tabularx}{\linewidth}{>{\bfseries}l X X}
\toprule
Bénéfice & Ce que le TinyML apporte & Nuance à garder en tête \\
\midrule
Réactivité      & Décision locale, sans aller-retour réseau & La latence est réduite, pas supprimée \\
Fonctionnement hors ligne & Analyse maintenue sans connexion & Seulement pour les fonctions conçues pour cela \\
Confidentialité & Les données brutes peuvent rester dans l'objet & Dépend des choix réels de transmission \\
Énergie         & Moins de transmissions radio & Le calcul local consomme aussi \\
Coût            & Moins de données envoyées et stockées dans le cloud & Coût de développement et de mise à jour \\
\bottomrule
\end{tabularx}
\caption{Synthèse des bénéfices du TinyML et de leurs limites.}
\label{tab:benefices}
\end{table}

\section{Domaines d'application et exemples d'utilisation}

\subsection{Industrie : surveillance des machines et maintenance}

Dans l'industrie, l'apprentissage automatique embarqué peut servir à surveiller
le fonctionnement des machines et à détecter des signes de dégradation. La
maintenance prédictive fait notamment partie des applications de l'IA en
périphérie : l'objectif est d'identifier des défauts avant qu'ils n'entraînent
des problèmes plus importants~\cite{zhang2021aiot}.

Le scénario illustratif présenté dans la partie~\ref{ch:fondements} relève de
ce domaine : un moteur surveillé par un capteur de vibrations, dont les mesures
seraient analysées sur un microcontrôleur afin de déclencher une alerte
lorsqu'un comportement inhabituel est détecté.

L'intérêt recherché est de fournir une information utile au personnel de
maintenance sans transmettre continuellement toutes les mesures à un serveur.
Il convient néanmoins de distinguer la détection d'anomalies de la prédiction
d'une panne : repérer un fonctionnement inhabituel ne permet pas nécessairement
de déterminer sa cause ou le temps restant avant une défaillance. La détection
constitue un premier niveau de surveillance ; la prédiction demande une analyse
et une validation supplémentaires.

\subsection{Maison connectée : reconnaissance de mots-clés et de sons}

La reconnaissance de mots-clés est une application représentative du TinyML ;
elle fait d'ailleurs partie des tâches de référence du banc d'essai MLPerf
Tiny~\cite{banbury2021mlperftiny}. Un modèle compact peut être chargé
d'identifier un nombre limité de commandes vocales sur un appareil disposant de
peu de mémoire et de puissance de calcul. Cette tâche ciblée se distingue d'une
conversation complète avec un assistant fondé sur un grand modèle de langage.

L'analyse locale présente ici un double intérêt : réduire les échanges avec un
serveur et permettre une réaction rapide. Elle peut aussi limiter la
transmission de données audio lorsque le système est conçu pour ne communiquer
que les résultats nécessaires.

La reconnaissance reste cependant sensible aux conditions d'utilisation. Le
bruit ambiant peut perturber la détection et provoquer des erreurs. Pour
obtenir un système robuste, il faut donc entraîner et tester le modèle avec des
données représentatives des environnements dans lesquels l'objet sera
réellement utilisé.
